

\section{Introduction\label{sec:introduction}}

This article compares and contrasts  methods developed in two different
literatures for inferring preferences (and sometimes also beliefs) of human decision makers 
from observations on their behavior over time: one developed by economists
that can be called ``structural econometrics'' (SE) and another developed by computer scientists 
in machine learning (ML) and artificial intelligence (AI) that is broadly under the general 
rubric called ``reinforcement learning'' (RL). Both SE and RL are huge literatures
and we will not attempt to survey them in any breadth. Instead, we will focus on a particular
subarea of SE known as {\it dynamic discrete choice\/} (DDC) which is most clearly
parallel to a subareas of RL called ``inverse reinforcement learning'' (IRL).\footnote{\footnotesize To space we will rely heavily
on abbreviations throughout this article, but refer readers to table \ref{tab:abbrevs} for a summary
of the various abbreviations and acronyms used throughout this article.}
Both literatures rely on the presumption that the agents whose preferences we are interested in recovering
are {\it rational\/} i.e. they take a sequence
of actions over time to maximize an internally perceived reward (or ``utility'') given their subjective beliefs of how
their current actions cause or at least influence potentially uncertain future outcomes that affect their rewards. The 
goals of DDC and IRL are identical:  to invert or infer the underlying objective or reward function (and sometimes also their beliefs) that is presumed to govern the behavior of agents 
from observations of the sequence of  states and actions of a sample of decision makers (DMs) over time.
The DMs could be  individuals, firms, or other intelligent agents (e.g. animals).
Beyond the intellectual interest in whether this type of inversion and inference is possible
what are the practical motivations for DDC and  IRL modeling?   

In economics, structural models are used to test and evaluate economic theories, and to make {\it counterfactual predictions\/} 
of the effect of changes in
economic policies or other changes in the environment on behavior and welfare 
(i.e. who is hurt and who is helped by the policy change?). Structural models are distinguished
from reduced-form models that provide statistical summaries of agents' behavior, but do not attempt to
infer their preferences and beliefs. Reduced-form models can make counterfactual predictions of a policy change when there
has been sufficient historical variation in policy to allow it to be included as an explanatory variable in 
an econometric model that predicts economic behavior. However reduced form models are unable to predict their impact on behavior for policies that have not changed in the past whereas 
structural models can produce accurate ``out-of-sample''
predictions of policy changes on behavior and welfare in the absence of past variation in the policy. 

For example, 
for decades Denmark has had one of the highest tax rates on new cars of any 
country in the world. Though there have been changes in the new car registration tax rate,
it has been to encourage more purchases of a relatively small share
of low emissions vehicles. \citet{GIMRS:2022} used a structural model of equilibrium ownership and trading of cars in Denmark
to predict the impact of reducing this ``registration tax''. The model predicts that by reducing the registration
tax and increasing the gas tax, it would be possible of Denmark to 1) improve the welfare of its citizens,
2) increase total tax revenue from cars (registration plus gas tax), and 3) reduce total $\mbox{CO}_2$ pollution.
It can be substantially faster and cheaper to use structural models to predict and evaluate the effects of 
various changes in economic policies than to actually carry them out and learning over time by trial and error. 
Auto producers learned this lesson decades ago when they developed finite element simulations
of car crashes that were sufficiently realistic that they were adopted as a faster and cheaper way to evaluate new car designs.
Structural models can be used to ``crash test'' new policy ideas in a similar way.

The practical motivation for inferring the DM's reward function in IRL  stems from the desire to produce
intelligent agents that can perform well in situations
with complex, poorly-defined, or hard to specify goals. They do this by inferring rewards
that motivate humans to perform  these same tasks,  using limited human {\it demonstration data.\/} 
ML and RL have had astounding success in training
machines to understand, predict, and generate text, artwork, video, spoken language, in addition to impressive strides in robotics.
This success depends not only on advances in computer power and
algorithms such as deep neural networks (DNN) that have the flexibility to make good predictions from a variety
of unstructured inputs, but also on the existence of large amounts of training data, {\it and\/} a well defined objective
function that the  algorithm  is trained to optimize. 

As \citet{RussellNorvig:2022} note, most of the existing ML methods use {\it supervised learning\/} where the algorithm
``learns by passively observing example input/output pairs provided by a `teacher'.'' (p. 840). Here the reward function is known:
to predict the choices of the teacher as well as possible according to some metric. However the labelled data provided a by teacher
is often quite limited, and this is where RL has proven so successful. When a reward function is known
RL algorithms can produce behavior that approximately maximizes the reward function 
without an explicit teacher. A good example is chess, where  the reward function is known (1 for a win, -1 for a loss
and 0 for a draw).  \citet{RussellNorvig:2022} note that supervised learning of chess is limited by the 
``relatively few examples (about $10^8$) compared to the space of all possible
chess positions (about $10^{40}$)'' whereas RL has been able to learn superhuman peformance 
via  {\it self-play,\/} i.e. repeatedly playing the algorithm against computerized opponents in
``offline training'' resulting in performance that exceeds the best human and computerized chess players 
in ``online'' real-time play, see e.g. \citet{Silver:2018}.   

Still, there are many challenging tasks in AI and robotics where it is unclear what reward function should be 
used by RL to train intelligent behavior.  \citet{christiano:2017} provide a concrete example:  
``suppose that we wanted to use reinforcement learning to train a robot to clean a table or scramble an egg. It's not clear how to construct a suitable reward function, which will need to be a function of the robot's sensors'' and the option of ``manually programming their behavior has become increasingly challenging and expensive.'' \citet{Osa:2018}. 
The IRL literature circumvents this by inferring a reward function from limited observations of human choices 
for a task we would like to automate. If it is possible to estimate 
a reward function that ``rationalizes'' human choices,  RL can use it 
to train the algorithm to optimize it much more rapidly and cheaply under a wider range of scenarios
that were encountered in the demonstration data set. The result is intelligent  behavior 
consistent with human preferences, including situations that are outside 
the limited set in the data it was trained on.\footnote{\footnotesize There is also a literature
on ``inverse optimal control'' (IOC) originated by \citet{kalman:1964} that is also related to structural estimation and IRL.  One of the early
practical successes in this literature, \citet{locomotion:2010}, was in robotics where IOC was used to ``generate natural overall locomotion trajectories from an initial rest position an orientation to a given target rest position and orientation.''
\citet{guidedcostlearning:2016} extend the IOC approach to use DNNs to provide more flexible approximations to the
underlying reward or cost function  that ``can learn complex, nonlinear cost representations $\ldots$ and can be applied to high-dimensional systems with unknown dynamics'' resulting in a method that ``outperforms prior IOC algorithms on a set of simulated benchmarks, and achieves good results on several real-world tasks.'' See \citet{azar:2020} for further connections between IOC and IRL.}

For example \citet{barnes:2024} used data on trips from Google Maps to ``provide routes that best reflect
travelers' latent preferences'' which can only be inferred from their ``physical behavior, which implicitly
trade-off factors including traffic conditions, distance, hills, safety, scenery, road conditions, etc.'' They showed
how to massively scale IRL to estimate a ``route preference function'' with 360 million parameters using  
a training dataset of 110 million trips sampled from Google Maps.  Using this reward function, they 
used  RL to train Google Maps to provide  route recommendations more consistent with
human preferences resulting ``in a policy that achieves a 16-24\% improvement
in route quality at a global scale, and to the best of our knowledge, represents the largest published study or IRL
algorithms in a real-world setting to date.''\footnote{\footnotesize While 110 million observations may not seem
``limited'', it is a tiny fraction of all trips navigated using Google Maps. \citet{Jeon:2021} cites
additional examples where ``defining a reward function beforehand is particularly challenging and IRL is simply
more pragmatic.'' These examples include robotic maniuplation, autonomous driving, and clinical motion analysis.
See \citet{zhao:2023} who used Adversarial IRL (AIRL) to estimate context-dependent
route preferencs using data on trips by taxi drivers in Shanghai.} 

The remainder of this article is organized as follows. Section \ref{sec:mdp} provides background on dynamic programming
(DP) and Markovian decision processes (MDP) which is the common ancestor of both the SE and RL literatures. We introduce
the key concepts of value and policy function and the Bellman equation which characterizes the solution for an
optimal dynamic policy, and introduce RL as a class of iterative stochastic algorithms for solving the MDP problem. 
Section \ref{sec:ddc} describes how the DP framework is adapted to inference of the
reward function in the DDC literature which the method of maximum likelihood. We also discuss the
{\it identification problem\/} that puts inherent limits on our ability to uniquely recover the true preferences
 and beliefs of a decision maker in the absences of strong additional assumptions on their functional form.
Section \ref{sec:irl} summarizes the literature on IRL and its relation to the DDC and section \ref{sec:conclusion} 
contrast the DDC and IRL literatures and discusses several unsolved problems.


\begin{table}[h!]
\centering
\caption{List of abbreviations and acronyms}
\begin{tabular}{|l|l|}
\hline
\textbf{Abbreviation} & \textbf{Full Name} \\
\hline
AI & Artificial Intelligence \\
AIRL & Adversarial Inverse Reinforcement Learning\\
AGI & Artificial General Intelligence \\
AL & Apprenticeship Learning \\
BC & Behavioral Cloning \\
BIRL & Bayesian Inverse Reinforcement Learning\\
CCP & Conditional Choice Probability \\
CCS & Conditional Choice Simulator \\
DDC & Dynamic Discrete Choice \\
DM  & Decision Maker (Agent) \\
DNN & Deep Neural Network \\
DP  & Dynamic Programming \\
DSE & Dynamic Structural Estimation \\
DU & Discounted Utility \\
EU   & Expected Utility\\
IOC & Inverse Optimal Control \\
IL & Imitation Learning  \\
IRL & Inverse Reinforcement Learning\\
MCMC & Markov Chain Monte Carlo \\
MDP & Markovian Decision Process \\
ML & Machine Learning \\
MPE & Markov Perfect Equilibrium \\
MPEC & Mathematical Programming with Equilibrium Constraints\\
NFXP & Nested Fixed Point \\
%PG & Policy Gradient \\
PI  & Policy Iteration \\
%PPO & Proximal Policy Optimization \\
RFE & Reduced Form Estimation \\
RL & Reinforcement Learning  \\
RTDP & Real Time Dynamic Programming \\
SA  & Successive Approximations \\
TD  & Temporal Difference\\
%TRPO & Trust Region Polcy Optimization \\
%XAI, XRL & Explainable AI, RL \\
\hline
\end{tabular}
\label{tab:abbrevs}
\end{table}
